%============================================================
% LITERATURE REVIEW
%============================================================

\section{Literature Review}

\subsection{Count Data and the Poisson Distribution}

Count outcomes occur frequently in medical and epidemiological
research. Examples include the number of hospital visits, disease
episodes, adverse events, and days with symptoms. The Poisson
distribution is one of the standard probability models for such
outcomes because it provides a simple representation of event counts
through a single rate parameter.

A key property of the Poisson distribution is equidispersion, meaning
that the mean and variance are equal. This assumption can be
restrictive in applications where individuals differ in their
underlying event rates. When the observed variance is substantially
larger than the mean, a standard Poisson model may therefore provide
a poor description of the data. The assessment of such assumptions
forms an important part of count-data modelling
\citep{MolenberghsVerbeke2017}.

Over-dispersion can arise for several reasons, including unobserved
heterogeneity between individuals, dependence between events, and
model misspecification. Alternative count-data models have therefore
been developed to accommodate departures from the standard Poisson
assumption, including negative binomial and COM-Poisson models
\citep{Sellers2012,Azevedo2023}.

For the present study, the main interest is not simply to allow for
over-dispersion but to investigate whether the excess variation in
headache frequency may reflect distinct underlying groups of
patients. This motivates the use of finite mixture models.


%============================================================
\subsection{Heterogeneity and Finite Mixture Models}

Finite mixture models provide a way of representing a population as
a combination of several unobserved component distributions. Instead
of assuming that all observations arise from the same probability
distribution, the population is assumed to contain a finite number
of latent groups, each with its own parameters and mixing proportion
\citep{McLachlanPeel2000}.

For count data, a Poisson finite mixture model assumes that the
observed counts arise from one of several Poisson distributions. The
latent component is not observed directly, but the probability that
an observation belongs to each component can be estimated from the
data.

This approach is useful when over-dispersion is thought to result
from differences between subpopulations. A mixture of Poisson
distributions can have a variance greater than its mean even though
each individual component follows a Poisson distribution. The
additional variation arises from differences between the component-
specific Poisson means.

Finite mixture models have been applied in a wide range of areas,
including medical research, where populations may contain individuals
with different underlying disease or risk profiles
\citep{McLachlanPeel2000}. In the context of headache frequency, this
provides a useful framework for investigating whether patients can be
described by different underlying levels of headache burden rather
than by a single population-average rate.


%============================================================
\subsection{Nonparametric Maximum Likelihood Estimation}

Before specifying a fixed number of mixture components, the
underlying mixing distribution can be investigated using
nonparametric maximum likelihood estimation (NPMLE). Rather than
assuming in advance that the mixing distribution has a particular
number of support points, NPMLE estimates the distribution directly
from the observed data.

The theoretical basis for NPMLE in mixture models was developed by
Lindsay, who examined the geometry of mixture likelihoods and the
properties of estimated mixing distributions
\citep{Lindsay1983,Lindsay1983b}.

NPMLE is useful in the present analysis because it provides an
exploratory assessment of whether the data support a homogeneous
population or several underlying support points. The gradient
function can then be examined to identify regions that may correspond
to candidate mixture components.

This provides a link between exploratory analysis and the subsequent
finite mixture models. Rather than choosing the number of components
only by repeatedly fitting increasingly complex models, the NPMLE
analysis provides information about the possible structure of the
underlying mixing distribution.


%============================================================
\subsection{Estimation of Finite Mixture Models}

Finite mixture models contain an unobserved component-membership
variable, which makes direct maximum likelihood estimation more
complicated. The expectation-maximization (EM) algorithm provides a
general approach for estimating models with this type of missing or
unobserved information \citep{McLachlanPeel2000}.

The EM algorithm alternates between two stages. In the expectation
step, the current parameter estimates are used to calculate the
posterior probability that each observation belongs to each mixture
component. In the maximization step, the parameters are updated using
these probabilities. The process is repeated until convergence.

The EM framework is particularly relevant to the present study because
component membership is not directly observed. The estimated
posterior probabilities provide both a means of estimating the model
parameters and a basis for classifying individual patients according
to their most likely latent headache-frequency group.


%============================================================
\subsection{Determining the Number of Components}

An important issue in finite mixture modelling is determining how many
components should be included. Adding components generally improves
the likelihood, but excessive complexity can result in overfitting.
Model selection therefore requires consideration of both goodness of
fit and model complexity \citep{McLachlanPeel2000}.

Information criteria provide a practical approach to this problem.
The Akaike Information Criterion (AIC), introduced by Akaike
\citep{Akaike1974}, balances model fit against the number of estimated
parameters. Smaller AIC values indicate a preferred trade-off
between fit and complexity.

The Bayesian Information Criterion (BIC), proposed by Schwarz
\citep{Schwarz1978}, applies a stronger penalty for model complexity
than AIC. It is therefore possible for AIC and BIC to select
different numbers of components when candidate models have similar
likelihoods but different numbers of parameters.

In finite mixture analysis, component selection should therefore not
depend on a single numerical criterion. Information criteria can be
considered alongside likelihood values, diagnostic plots, estimated
mixing proportions, and the interpretability of the resulting
components \citep{McLachlanPeel2000}.

This consideration is particularly relevant to the present study,
where the two- and three-component models provide different results
under AIC and BIC.


%============================================================
\subsection{Mixture Models and Headache Frequency}

Headache frequency is naturally represented as a count outcome when
the outcome is defined as the number of headache days during a fixed
period. The clinical data used in this study originate from a
randomized trial investigating acupuncture for chronic headache in
primary care \citep{Vickers2004}.

A conventional Poisson model provides an initial framework for such
data, but it assumes that patients share a common underlying rate
after accounting for included covariates. This assumption may be
unrealistic when there are substantial differences in headache
frequency between patients.

Patients with relatively few headache days may have a different
underlying rate from patients who experience headaches on most days
of the observation period. A finite mixture model provides a way to
represent this possibility without requiring the latent groups to be
defined in advance.

The latent components are estimated from the observed distribution
and can subsequently be examined for their statistical and clinical
interpretability. The present study therefore uses Poisson finite
mixture modelling to investigate whether the heterogeneity in
baseline headache frequency can be represented by a small number of
latent groups.


%============================================================
\subsection{Summary of the Literature}

The literature provides a statistical basis for the modelling
strategy used in this study. The Poisson distribution provides a
natural starting point for count outcomes, but its equidispersion
assumption can be restrictive when substantial heterogeneity is
present. Finite mixture models provide one approach to representing
this heterogeneity through latent subpopulations
\citep{MolenberghsVerbeke2017,McLachlanPeel2000}.

NPMLE provides a flexible approach for exploring the underlying
mixing distribution, while the EM algorithm provides a practical
method for estimating finite mixture parameters
\citep{Lindsay1983,McLachlanPeel2000}. AIC and BIC provide
complementary criteria for comparing models with different numbers
of components \citep{Akaike1974,Schwarz1978}.

Together, these methods provide the framework used in this study to
investigate the distribution of headache frequency and to determine
whether a finite mixture model provides a more informative
description than a single Poisson distribution.